% Part: first-order-logic
% Chapter: introduction
% Section: soundness-completeness

\documentclass[../../../include/open-logic-section]{subfiles}

\begin{document}

\olfileid{fol}{int}{scp}

\olsection{বিশুদ্ধতা ও পূর্ণতা}

প্রথম-ক্রমের যুক্তিবিদ্যার জন্য !!{derivation} পদ্ধতিও প্রবর্তন করব।
যুক্তিবিদেরা অনেক !!{derivation} পদ্ধতি নির্মাণ করেছেন, তবে তারা সবাই
!!{sentence}s-এর মধ্যে একই !!{derivability} সম্পর্ক সংজ্ঞায়িত করে।
সুনির্দিষ্টভাবে সংজ্ঞায়িত কোনো ধরনের !!a{derivation} থাকলে আমরা বলি,
$\Gamma$,~$!A$-কে \emph{!!{derive}s} করে, $\Gamma \Proves !A$।
!!^{derivation}s সব সময় সংকেতের সসীম বিন্যাস—কখনো !!{sentence}s-এর
তালিকা, কখনো আরও জটিল কোনো গঠন। !!{derivation} পদ্ধতির উদ্দেশ্য হলো,
কোনো সেট~$\Gamma$ থেকে কোনো !!a{sentence} অনুসৃত হয় কি না তা নির্ধারণের
উপকরণ দেওয়া। সেই উদ্দেশ্য পূরণ করতে হলে $\Gamma \Entails !A$ হবে যদি
এবং কেবল যদি $\Gamma \Proves !A$।

$\Gamma \Proves !A$ হলেও $\Gamma \Entails !A$ না হলে আমাদের
!!{derivation} পদ্ধতিটি অতিরিক্ত শক্তিশালী—সেটি প্রয়োজনের বেশি প্রমাণ
করে। $\Gamma \Proves !A$ হলে $\Gamma \Entails !A$—এই ধর্মকে
\emph{বিশুদ্ধতা} বলা হয়; যেকোনো ভালো !!{derivation} পদ্ধতির জন্য এটি
ন্যূনতম শর্ত। অন্যদিকে, $\Gamma \Entails !A$ হলেও $\Gamma \Proves !A$
না হলে আমাদের !!{derivation} পদ্ধতিটি অতিরিক্ত দুর্বল—সেটি যথেষ্ট প্রমাণ
করে না। $\Gamma \Entails !A$ হলে $\Gamma \Proves !A$—এই ধর্মকে
\emph{পূর্ণতা} বলা হয়। সাধারণত বিশুদ্ধতা প্রমাণ করা তুলনামূলকভাবে সহজ
(আরোহীভাবে সংজ্ঞায়িত !!{derivation}s-এর গঠনের উপর আরোহ করে)। পূর্ণতা
প্রমাণ করা কঠিনতর।

বিশুদ্ধতা ও পূর্ণতার কয়েকটি গুরুত্বপূর্ণ পরিণাম আছে। !!{sentence}s-এর
কোনো সেট~$\Gamma$ যদি একটি স্ববিরোধ (যেমন $!A \land \lnot !A$)
!!{derive}s করে, তবে সেটিকে \emph{অসঙ্গত} বলা হয়। অসঙ্গত~$\Gamma$-র কোনো
মডেল থাকতে পারে না; সেগুলি অপরিতৃপ্তিযোগ্য। পূর্ণতা থেকে বিপরীত দাবিটি
অনুসৃত হয়: অসঙ্গত নয় এমন—অর্থাৎ, আমরা যাকে \emph{সঙ্গত} বলব—প্রতিটি
$\Gamma$-র একটি মডেল আছে। বস্তুত, এটি পূর্ণতার সমতুল্য এবং পূর্ণতার যে
রূপটি আমরা আসলে প্রমাণ করব সেটিই এটি। ফলটি গভীর এবং হয়তো বিস্ময়কর:
$\Gamma$ থেকে $!A \land \lnot !A$ প্রমাণ করা যায় না—শুধু এই তথ্যই
$\Gamma$-র বর্ণনার মতো !!a{structure}-এর অস্তিত্বের নিশ্চয়তা দেয়।
অতএব পূর্ণতা এই প্রশ্নের উত্তর দেয়: !!{sentence}s-এর কোন সেটগুলির মডেল
আছে? উত্তর: ঠিক সব সঙ্গত সেটেরই মডেল আছে।

বিশুদ্ধতা ও পূর্ণতা উপপাদ্যের দুটি গুরুত্বপূর্ণ পরিণাম হলো সংহতি উপপাদ্য
ও লোয়েনহাইম--স্কোলেম উপপাদ্য। এগুলি মডেল তত্ত্বের গুরুত্বপূর্ণ ফল এবং
এদের সাহায্যে বহু আকর্ষণীয় ফল প্রতিষ্ঠা করা যায়। দুটির উল্লেখ আমরা
ইতিমধ্যেই করেছি: প্রথম-ক্রমের যুক্তিবিদ্যায় কোনো !!a{structure}-এর
!!{domain} সসীম অথবা !!{nonenumerable}—এ কথা প্রকাশ করা যায় না।

ঐতিহাসিকভাবে এই সব কিছু—প্রথম-ক্রমের যুক্তিবিদ্যার সংকেতবিন্যাস ও
অর্থতত্ত্ব কীভাবে সংজ্ঞায়িত করতে হয়, ভালো !!{derivation} পদ্ধতি কীভাবে
সংজ্ঞায়িত করতে হয়, সেগুলি যে বিশুদ্ধ ও পূর্ণ তা কীভাবে প্রমাণ করতে হয়,
প্রথম-ক্রমের ভাষায় কী প্রকাশ করা যায় ও কী যায় না সে বিষয়ে স্পষ্ট
ধারণা পাওয়া—বুঝে উঠতে ও ঠিক করতে দীর্ঘ সময় লেগেছে। এখন আমরা জানি
কাজগুলি কীভাবে করতে হয়, কিন্তু সব খুঁটিনাটি পার হওয়া এখনও বিভ্রান্তিকর
ও ক্লান্তিকর হতে পারে। তবু তা গুরুত্বপূর্ণও, কারণ প্রথম-ক্রমের
যুক্তিবিদ্যার বিধিবদ্ধ ভাষার জন্য এখানে নির্মিত পদ্ধতিগুলি যুক্তিবিদ্যা,
কম্পিউটার বিজ্ঞান ও ভাষাবিজ্ঞানের সর্বত্র প্রয়োগ করা হয়। তাই
খুঁটিনাটিগুলি নিয়ে কাজ করলে শেষ পর্যন্ত তার সুফল পাওয়া যায়।

\end{document}
